% Libro 1 · Mundo 6 · Teoría de números (app: pu-numeros)
\mundo{Teoría de números}{pu-numeros}{Los números enteros esconden una estructura: cómo se descomponen en primos, cuándo
uno cabe exacto en otro, qué residuo deja una división. El examen la usa para preguntas que parecen imposibles sin
calculadora («la última cifra de $3^{2027}$») y que salen en tres renglones con la idea correcta.}{%
\item \textbf{6.1} Primos y divisibilidad
\item \textbf{6.2} Máximo común divisor y mínimo común múltiplo
\item \textbf{6.3} Paridad, residuos y última cifra}

\leccion{pu-nu-1}{Primos y divisibilidad}{Juegos del árbol de factores y la criba.}

\begin{reto}
¿Cuántos divisores positivos tiene el número 360?
\opciones{12}{18}{24}{36}
\end{reto}

\subsection{La idea: los primos son los ladrillos}
Un número \clave{primo} tiene exactamente dos divisores: 1 y él mismo (2, 3, 5, 7, 11, 13, 17, 19, 23, 29, …). El 1 no
es primo. Todo entero mayor que 1 se escribe de \emph{una sola manera} como producto de primos: es su «código genético».

\figura{factores}{El árbol de factores de 360: en verde, los primos.}

\subsubsection{Contar divisores}
Un divisor de $360 = 2^3 \cdot 3^2 \cdot 5$ se arma eligiendo cuántos 2 lleva (0, 1, 2 o 3: cuatro opciones), cuántos 3
(0, 1 o 2: tres) y cuántos 5 (0 o 1: dos). Por el principio multiplicativo hay $4 \cdot 3 \cdot 2 = 24$ divisores.

\subsubsection{Criterios de divisibilidad}
\begin{center}\small
\begin{tabular}{@{}cl@{}}
\toprule
\textbf{Entre} & \textbf{Se divide si…}\\\midrule
2 & la última cifra es par\\
3 & la suma de las cifras es múltiplo de 3\\
4 & las dos últimas cifras forman un múltiplo de 4\\
5 & termina en 0 o 5\\
6 & es divisible entre 2 y entre 3\\
8 & las tres últimas cifras forman un múltiplo de 8\\
9 & la suma de las cifras es múltiplo de 9\\
11 & la suma alternada de cifras ($+ - + \ldots$) es múltiplo de 11\\
\bottomrule
\end{tabular}
\end{center}

\subsubsection{¿Es primo?}
Para saber si $n$ es primo basta probar los primos hasta $\sqrt n$: si $n$ tuviera un factor mayor que $\sqrt n$,
tendría otro menor. ¿Es primo 221? $\sqrt{221} < 15$; se prueba 2, 3, 5, 7, 11, 13: $221 = 13 \cdot 17$. No.

\begin{regla}
\begin{itemize}
  \item Si $n = p^a q^b r^c$, el número de divisores positivos es $(a + 1)(b + 1)(c + 1)$.
  \item $n$ es un cuadrado perfecto si y solo si todos los exponentes son pares (y entonces tiene un número impar de
  divisores).
  \item Si $d$ divide a $a$ y a $b$, divide a $a + b$, a $a - b$ y a cualquier combinación $ma + nb$.
\end{itemize}
\end{regla}

\subsection{Ejemplos resueltos}
\begin{ejemplo}[Completar un criterio]
¿Qué cifra $d$ hace que $5\,d\,4\,2$ sea divisible entre 6?
\solucion Termina en 2, es par. Para el 3: $5 + d + 4 + 2 = 11 + d$ múltiplo de 3: $d = 1, 4, 7$. Hay tres cifras posibles.
\end{ejemplo}

\begin{ejemplo}[El menor múltiplo cuadrado]
¿Por qué número mínimo hay que multiplicar 180 para obtener un cuadrado perfecto?
\solucion $180 = 2^2 \cdot 3^2 \cdot 5$. Al 5 le falta otro 5: se multiplica por 5 y se obtiene $900 = 30^2$.
\end{ejemplo}

\begin{ejemplo}[Divisibilidad entre 11]
¿Es 918\,082 divisible entre 11?
\solucion Suma alternada desde la derecha: $2 - 8 + 0 - 8 + 1 - 9 = -22$, múltiplo de 11. Sí.
\end{ejemplo}

\begin{ejemplo}[Factores en un factorial]
¿Cuántas veces aparece el factor 2 en $10! = 1 \cdot 2 \cdots 10$?
\solucion Pares hasta 10: cinco (aportan un 2 cada uno). Múltiplos de 4: dos (aportan otro). De 8: uno (otro más).
Total $5 + 2 + 1 = 8$.
\end{ejemplo}

\begin{trampa}
Creer que 1 es primo, o que todos los impares son primos ($9$, $15$, $21$, $27$, $51 = 3 \cdot 17$, $91 = 7 \cdot 13$
no lo son).
\end{trampa}

\taller
\piso{Calienta}{descomponer}
\begin{ejercicios}[2]
\item Descompón 504. \sol{$2^3 \cdot 3^2 \cdot 7$.}
\item ¿Cuántos divisores tiene 72? \sol{$2^3 3^2$: $4 \cdot 3 = 12$.}
\item ¿Es primo 187? \sol{No: $11 \cdot 17$.}
\item ¿Es 7\,416 divisible entre 8? \sol{Sí: $416 = 8 \cdot 52$.}
\end{ejercicios}
\piso{Entrena}{problemas}
\begin{ejercicios}
\item ¿Cuántos números del 1 al 100 son divisibles entre 3 o entre 5? \sol{$33 + 20 - 6 = 47$.}
\item ¿Cuál es el menor número con exactamente 6 divisores? \sol{$12 = 2^2 \cdot 3$ ($3 \cdot 2 = 6$).}
\item ¿Cuántos divisores impares tiene 360? \sol{Sin el 2: $3^2 \cdot 5$ tiene $3 \cdot 2 = 6$.}
\end{ejercicios}
\piso{Como en el examen}{opción múltiple}
\begin{ejercicios}
\item ¿Cuál de estos números es primo?
\opciones{91}{87}{97}{111}
\sol{C. $91 = 7 \cdot 13$, $87 = 3 \cdot 29$, $111 = 3 \cdot 37$.}
\item El número $\overline{4a7b}$ es divisible entre 5 y entre 9 y $b \ne 0$. Entonces $a$ es
\opciones{0}{2}{1}{9}
\sol{B. Termina en 0 o 5 y $b \ne 0$, así que $b = 5$. La suma $4 + a + 7 + 5 = 16 + a$ debe ser múltiplo de 9: $a = 2$.}
\item ¿Cuántos divisores tiene $2^4 \cdot 3^3$?
\opciones{7}{12}{20}{24}
\sol{C. $5 \cdot 4 = 20$.}
\end{ejercicios}

\begin{retoresuelto}
$360 = 2^3 \cdot 3^2 \cdot 5^1$. Divisores: $(3 + 1)(2 + 1)(1 + 1) = 24$. \textbf{Opción C.}
\end{retoresuelto}
\bitacora

\leccion{pu-nu-2}{MCM y MCD}{Juegos de las luces que coinciden y la baldosa más grande.}

\begin{reto}
Tres luces parpadean cada 12, 18 y 30 segundos. Si parpadean juntas a las 8:00:00, ¿a qué hora vuelven a coincidir?
\opciones{8:01:00}{8:03:00}{8:06:00}{8:09:00}
\end{reto}

\subsection{La idea: ¿cuándo coinciden? ¿en qué se parten?}
Dos preguntas, dos herramientas:
\begin{itemize}
  \item «¿Cuándo vuelven a coincidir?», «¿cuál es el menor número que…?» $\to$ \clave{mínimo común múltiplo} (MCM): el
  menor número que es múltiplo de todos.
  \item «¿En cuántos trozos iguales, lo más grandes posible?», «¿cuál es la mayor baldosa?» $\to$ \clave{máximo común
  divisor} (MCD): el mayor número que divide a todos.
\end{itemize}

Con factores primos: el MCD toma los primos \emph{comunes} con el \emph{menor} exponente; el MCM toma \emph{todos} los
primos con el \emph{mayor} exponente.

\begin{center}
\begin{tabular}{rl}
$12 =$ & $2^2 \cdot 3$\\
$18 =$ & $2 \cdot 3^2$\\
$30 =$ & $2 \cdot 3 \cdot 5$\\\midrule
MCD $=$ & $2 \cdot 3 = 6$\\
MCM $=$ & $2^2 \cdot 3^2 \cdot 5 = 180$\\
\end{tabular}
\end{center}

\figura{mcm}{Dos ritmos distintos que coinciden cada cierto tiempo.}

\begin{regla}
\begin{itemize}
  \item Para dos números: $\text{MCD}(a, b) \cdot \text{MCM}(a, b) = a \cdot b$.
  \item Si $a$ y $b$ no tienen factores comunes ($\text{MCD} = 1$), su MCM es $ab$.
  \item Algoritmo de Euclides: $\text{MCD}(a, b) = \text{MCD}(b, a - b)$; restando (o dividiendo) se llega rápido.
\end{itemize}
\end{regla}

\subsection{Ejemplos resueltos}
\begin{ejemplo}[Baldosas]
Un piso mide 360 cm por 240 cm. ¿Cuál es la mayor baldosa cuadrada que lo cubre sin cortar, y cuántas se necesitan?
\solucion $\text{MCD}(360, 240) = 120$. Baldosas de 120 cm: $3 \cdot 2 = 6$.
\end{ejemplo}

\begin{ejemplo}[Un residuo común]
¿Cuál es el menor número mayor que 1 que al dividirlo entre 4, 6 y 9 deja residuo 1?
\solucion El número menos 1 es múltiplo de 4, 6 y 9: de $\text{MCM} = 36$. El número es 37.
\end{ejemplo}

\begin{ejemplo}[Euclides]
Calcula $\text{MCD}(391, 299)$.
\solucion $391 - 299 = 92$; $\text{MCD}(299, 92)$: $299 = 3 \cdot 92 + 23$; $\text{MCD}(92, 23) = 23$ porque $92 = 4
\cdot 23$. Es 23.
\end{ejemplo}

\begin{ejemplo}[Producto y MCD]
Dos números tienen MCD 6 y MCM 180. Si uno es 36, ¿cuál es el otro?
\solucion $6 \cdot 180 = 36 \cdot x$: $x = 30$.
\end{ejemplo}

\begin{trampa}
Usar el MCD cuando se pregunta por coincidencias en el tiempo, o el MCM cuando se pregunta por partir en trozos. Pregunta
de control: ¿el resultado debe ser \emph{mayor} que los datos (MCM) o \emph{menor} (MCD)?
\end{trampa}

\taller
\piso{Calienta}{calcular}
\begin{ejercicios}[2]
\item MCD y MCM de 24 y 36. \sol{12 y 72.}
\item MCM de 8, 12 y 15. \sol{120.}
\item MCD de 84, 126 y 210. \sol{42.}
\item MCM de 7 y 9. \sol{63.}
\end{ejercicios}
\piso{Entrena}{problemas}
\begin{ejercicios}
\item Dos buses salen de la terminal cada 15 y 20 minutos. Si salen juntos a las 6:00, ¿cuántas veces salen juntos hasta las 12:00 (incluidas)? \sol{Cada 60 min: 6:00, 7:00, …, 12:00: 7 veces.}
\item Se tienen 48 rosas y 72 claveles para armar ramos iguales con todas las flores, el mayor número posible de ramos. ¿Cuántos ramos y de qué composición? \sol{MCD 24: 24 ramos de 2 rosas y 3 claveles.}
\item ¿Cuál es el menor número de 3 cifras divisible entre 6, 8 y 10? \sol{MCM 120: el menor es 120.}
\end{ejercicios}
\piso{Como en el examen}{opción múltiple}
\begin{ejercicios}
\item Tres atletas dan una vuelta a la pista en 60, 72 y 90 segundos. Salen juntos. Vuelven a coincidir en la salida a los
\opciones{6 min}{12 min}{18 min}{3 min}
\sol{A. MCM $= 360$ s.}
\item Se quiere cortar dos varillas de 180 cm y 252 cm en trozos iguales, lo más largos posible. ¿Cuántos trozos salen en total?
\opciones{12}{7}{36}{24}
\sol{A. MCD 36: $5 + 7 = 12$ trozos.}
\item Si $\text{MCD}(a, 40) = 8$ y $\text{MCM}(a, 40) = 120$, entonces $a$ es
\opciones{24}{48}{15}{12}
\sol{A. $8 \cdot 120 / 40 = 24$.}
\end{ejercicios}

\begin{retoresuelto}
$\text{MCM}(12, 18, 30) = 180$ segundos $= 3$ minutos. Vuelven a coincidir a las 8:03:00. \textbf{Opción B.}
\end{retoresuelto}
\bitacora

\leccion{pu-nu-3}{Paridad, residuos y última cifra}{Juegos del reloj de residuos y la cifra que da vueltas.}

\begin{reto}
¿Cuál es la última cifra de $3^{2027}$?
\opciones{1}{3}{7}{9}
\end{reto}

\subsection{La idea: los residuos dan vueltas como un reloj}
Al dividir entre 4, los residuos solo pueden ser 0, 1, 2 o 3, y se repiten en ciclo: 0, 1, 2, 3, 0, 1, 2, 3, … Esto es
como un reloj de 4 horas. Muchas preguntas sobre números enormes solo dependen de «en qué hora del reloj» están.

\subsubsection{Últimas cifras de potencias}
La última cifra de un producto depende solo de las últimas cifras de los factores. Por eso las potencias forman ciclos:

\begin{center}\small
\begin{tabular}{@{}cl@{\qquad}cl@{}}
\toprule
\textbf{Base} & \textbf{Ciclo de la última cifra} & \textbf{Base} & \textbf{Ciclo}\\\midrule
2 & 2, 4, 8, 6 & 7 & 7, 9, 3, 1\\
3 & 3, 9, 7, 1 & 8 & 8, 4, 2, 6\\
4 & 4, 6 & 9 & 9, 1\\
0, 1, 5, 6 & siempre la misma & & \\\bottomrule
\end{tabular}
\end{center}

Para saber la última cifra de $3^{2027}$ basta saber en qué lugar del ciclo de 4 cae 2027: el residuo de $2027 \div 4$
es 3, así que es el tercero del ciclo.

\subsubsection{Paridad}
Par $+$ par $=$ par, impar $+$ impar $=$ par, par $+$ impar $=$ impar; un producto es impar solo si todos los factores
son impares. La paridad resuelve rápido preguntas de «¿es posible…?»: la suma de cinco números impares nunca es par.

\figura{relojciclo}{Las últimas cifras de las potencias de 3 dan vueltas.}

\begin{regla}
\begin{itemize}
  \item Para residuos de sumas y productos, se puede trabajar solo con los residuos: el residuo de $a \cdot b$ entre $n$
  es el residuo de (residuo de $a$) $\cdot$ (residuo de $b$).
  \item Ciclos: si un patrón se repite cada $k$ pasos, el paso $N$ es igual al paso $r$, donde $r$ es el residuo de
  $N \div k$ (si $r = 0$, es el último del ciclo).
  \item Días de la semana: avanzar $N$ días es avanzar $N \bmod 7$.
\end{itemize}
\end{regla}

\subsection{Ejemplos resueltos}
\begin{ejemplo}[Día de la semana]
Si hoy es martes, ¿qué día será dentro de 100 días?
\solucion $100 = 14 \cdot 7 + 2$: dos días después del martes, jueves.
\end{ejemplo}

\begin{ejemplo}[Residuo de una potencia]
¿Qué residuo deja $2^{50}$ al dividirlo entre 7?
\solucion Residuos de $2^1, 2^2, 2^3, \ldots$ entre 7: 2, 4, 1, 2, 4, 1, … ciclo de 3. $50 = 16 \cdot 3 + 2$: residuo 4.
\end{ejemplo}

\begin{ejemplo}[Última cifra de una suma]
¿Cuál es la última cifra de $7^{2026} + 8^{2026}$?
\solucion $2026 = 506 \cdot 4 + 2$. Para el 7, el segundo del ciclo: 9. Para el 8: 4. La suma termina en $9 + 4 = 13$: 3.
\end{ejemplo}

\begin{ejemplo}[Paridad]
¿Se pueden repartir los números del 1 al 10 en dos grupos con la misma suma?
\solucion La suma total es 55, impar: no se puede partir en dos mitades enteras iguales. No.
\end{ejemplo}

\begin{trampa}
Cuando el residuo es 0, la cifra es la \emph{última} del ciclo, no la primera: $3^{2028}$ termina en 1, no en 3.
\end{trampa}

\taller
\piso{Calienta}{ciclos}
\begin{ejercicios}[2]
\item Última cifra de $2^{100}$. \sol{$100 \bmod 4 = 0$: la cuarta, 6.}
\item Última cifra de $9^{55}$. \sol{Exponente impar: 9.}
\item Residuo de $1\,000$ entre 7. \sol{$1\,000 = 142 \cdot 7 + 6$: 6.}
\item Si hoy es sábado, ¿qué día fue hace 30 días? \sol{$30 = 4 \cdot 7 + 2$: dos días antes, jueves.}
\end{ejercicios}
\piso{Entrena}{problemas}
\begin{ejercicios}
\item ¿Qué letra ocupa el lugar 2026 en ZUKAZUKAZUKA…? \sol{$2026 \bmod 4 = 2$: U.}
\item ¿Cuál es el residuo de $13^{40}$ entre 12? \sol{$13$ deja residuo 1; $1^{40} = 1$.}
\item La suma de 4 números naturales consecutivos, ¿puede ser 2026? \sol{$4n + 6 = 2026$ da $n = 505$: sí (505 a 508).}
\end{ejercicios}
\piso{Como en el examen}{opción múltiple}
\begin{ejercicios}
\item La última cifra de $2^{2026}$ es
\opciones{2}{4}{6}{8}
\sol{B. $2026 \bmod 4 = 2$.}
\item El 1 de enero de cierto año fue lunes. Si el año no es bisiesto, el 31 de diciembre fue
\opciones{lunes}{martes}{domingo}{sábado}
\sol{A. Del 1 de enero al 31 de diciembre hay 364 días, múltiplo de 7.}
\item Si $n$ es impar, ¿cuál expresión es siempre par?
\opciones{$n^2$}{$n^2 + n$}{$3n$}{$n^2 + 2$}
\sol{B. $n(n + 1)$: producto de consecutivos.}
\end{ejercicios}

\begin{retoresuelto}
Ciclo de 4: 3, 9, 7, 1. $2027 = 506 \cdot 4 + 3$: el tercero del ciclo, 7. \textbf{Opción C.}
\end{retoresuelto}
\bitacora

\evaluacion
\begin{ejercicios}
\item ¿Cuántos divisores positivos tiene 200?
\opciones{8}{10}{12}{6}
\sol{C. $2^3 5^2$: $4 \cdot 3$. (6.1)}
\item ¿Qué cifra debe ir en $\overline{37\square 4}$ para que el número sea divisible entre 9?
\opciones{4}{5}{3}{0}
\sol{A. $3 + 7 + 4 = 14$; con 4 suma 18. (6.1)}
\item El menor número por el que hay que multiplicar 75 para obtener un cuadrado perfecto es
\opciones{3}{5}{15}{25}
\sol{A. $75 = 3 \cdot 5^2$. (6.1)}
\item El MCD de 150 y 225 es
\opciones{25}{75}{15}{450}
\sol{B. $2 \cdot 3 \cdot 5^2$ y $3^2 5^2$. (6.2)}
\item Un número entre 50 y 100 deja residuo 2 al dividirlo entre 3, entre 4 y entre 5. El número es
\opciones{60}{62}{64}{92}
\sol{B. Es $60 + 2$. (6.2)}
\item Dos engranajes de 24 y 36 dientes están encajados. ¿Cuántas vueltas da el pequeño hasta que vuelven a la posición inicial?
\opciones{2}{3}{6}{12}
\sol{B. MCM 72 dientes: el de 24 da 3 vueltas. (6.2)}
\item La última cifra de $7^{2026}$ es
\opciones{7}{9}{3}{1}
\sol{B. $2026 \bmod 4 = 2$. (6.3)}
\item ¿Qué residuo deja $5^{100}$ al dividirlo entre 4?
\opciones{0}{1}{2}{3}
\sol{B. 5 deja residuo 1. (6.3)}
\item Si hoy es miércoles, dentro de 1\,000 días será
\opciones{martes}{miércoles}{jueves}{domingo}
\sol{A. $1\,000 \bmod 7 = 6$: seis días después del miércoles. (6.3)}
\item La suma $1 + 2 + 3 + \cdots + 99$ es
\opciones{par y múltiplo de 3}{impar}{par y no múltiplo de 3}{primo}
\sol{A. $\tfrac{99 \cdot 100}{2} = 4\,950 = 2 \cdot 3^2 \cdot 5^2 \cdot 11$. (6.3)}
\end{ejercicios}
